\section{Introduction}\label{sec:introduction}

Let $K\subset\R^n$ be an origin-symmetric convex body: a compact convex
set with nonempty interior and $K=-K$. Its polar is
\[
 K^\circ=\{p\in\R^n:\ip{q}{p}\le1\text{ for every }q\in K\}.
\]
We place $K$ in the position coordinates and $K^\circ$ in the conjugate
momentum coordinates of $(\R^n_q\times\R^n_p,\omega_0)$, where
\[
 \omega_0=\sum_{j=1}^n dq_j\wedge dp_j,
 \qquad U_K=\interior K\times\interior K^\circ.
\]
For $c>0$, write
\[
 B^{2n}(c)=\{(q,p):\pi(|q|^2+|p|^2)<c\}.
\]
The Gromov width $c_G(U)$ of an open set $U\subset\R^{2n}$ is the supremum
of the capacities $c$ for which there is a smooth embedding
$e:B^{2n}(c)\to U$ satisfying $e^*\omega_0=\omega_0$.

\begin{theorem}\label{thm:main}
For every integer $n\ge2$ and every origin-symmetric convex body
$K\subset\R^n$,
\[
 c_G(U_K)=4.
\]
More precisely, for every $0<c<4$ there is a smooth symplectic embedding
$B^{2n}(c)\hookrightarrow U_K$.
\end{theorem}

The theorem imposes no smoothness, strict convexity, unconditional
symmetry, or product structure on $K$. It concerns every capacity
strictly below $4$; no embedding at the supremum is needed.
Symplectic embeddings preserve volume, and
$\vol_{2n}(B^{2n}(c))=c^n/n!$. Consequently Theorem~\ref{thm:main}
gives
\[
 \vol_n(K)\vol_n(K^\circ)\ge\frac{4^n}{n!}.
\]
Thus the theorem resolves the symmetric Mahler conjecture positively;
the complete deduction, including the one-dimensional case, appears in
Corollary~\ref{cor:mahler}. We make no claim here about the classification
of equality cases.

For $n\ge3$, the width theorem also transports finite ball packings
satisfying strict volume and pairwise-capacity inequalities into $U_K$; see
Corollary~\ref{cor:polar-packings}.

\subsection*{Background and prior constructions}

The symmetric volume-product problem goes back to Mahler's work on
transference in the geometry of numbers \cite{Mahler1939}.
Earlier sharp volume inequalities cover the plane
\cite[Theorem~13]{BoroczkyMakaiMeyerReisner2013}, unconditional bodies
(those invariant under coordinate sign changes)
\cite[Sections~1--2]{Reisner1987}, and dimension three
\cite{IriyehShibata}. Bourgain and Milman's asymptotic reverse
Santal\'o inequality gives a lower bound $a^n/n!$ for some $a>0$
independent of dimension \cite{BourgainMilman}; the sharp target is
$4^n/n!$. The companion paper \cite{SymmetricCompanion2026} discusses
the wider convex-geometric history and the equality problem. Our focus
is the stronger geometric assertion that the polar product contains
large symplectic balls.

Gromov's nonsqueezing theorem established that ball embeddings detect
symplectic restrictions beyond volume preservation \cite{Gromov1985}.
Viterbo later proposed a sharp volume--capacity inequality for convex
domains \cite{Viterbo2000}. Artstein-Avidan, Karasev, and Ostrover
connected this proposed inequality to the symmetric Mahler conjecture
and proved
\[
 c_{\mathrm{HZ}}(K\times K^\circ)=4
\]
for every origin-symmetric convex body
\cite[Theorems~1.6--1.7 and Section~4]{AAKO2014}.
Here $c_{\mathrm{HZ}}$ is the Hofer--Zehnder capacity, normalized to
have value $\pi R^2$ on a ball of radius $R$
\cite{HoferZehnder1994}. This computation implies the upper bound for
Gromov width. We use the supporting-cylinder argument of
\cite[Remark~4.2]{AAKO2014} and nonsqueezing to prove that bound directly.
The matching lower bound requires actual ball embeddings; it does not
follow from the Hofer--Zehnder capacity computation.

Earlier constructions establish this lower bound for particular
families. Karasev proves it for $\ell_p$ unit balls, $1<p<\infty$,
\cite[Proposition~3.1]{Karasev2021}, using coordinatewise area-preserving
constructions related to those of Latschev, McDuff, and Schlenk
\cite{LatschevMcDuffSchlenk2013}. Ramos and Sepe identify the interior of
the cube--crosspolytope product with the open ball of capacity $4$
\cite[Theorem~7]{RamosSepe2019}. Vicente obtains capacity equalities
for functional-dual products defined by Young functions and lower
bounds for ordinary polar products
\cite[Theorems~1.3 and~1.5]{Vicente2025}.
Functional duality and ordinary polarity give different partners, so
these results require separate hypotheses. Theorem~\ref{thm:main}
answers positively the unrestricted symmetric-polar-product width
question formulated in \cite{VicenteQuestion2026}.

The unrestricted Viterbo volume--capacity conjecture was disproved by
Haim-Kislev and Ostrover \cite{HKO2025}. Their four-dimensional
counterexample uses a regular pentagon and a rotated partner
\cite[Theorem~1.3 and Proposition~1.4]{HKO2025}; the pentagon is not
centrally symmetric. The theorem here concerns the particular
Lagrangian product of a symmetric body and its polar, and makes no
assertion for arbitrary convex domains in phase space.

The planar coordinate is Gross's uniform-distribution map in the
conformal Skorokhod construction, rotated in source and target
\cite[Section~3.2]{Gross2019}. The same lens and high powers of its
inverse coordinate appear in the complete companion paper
\cite[Section~3 and proof of Lemma~5.1]{SymmetricCompanion2026}.
Every lens estimate needed here, including the estimate uniform up to
the tips, is proved locally.
Neither the volume-product theorem nor the equality classification of
the companion is used in the proof.

The holomorphic-zero argument belongs to the broader use of logarithmic
poles to measure local positivity, as in Demailly
\cite[Section~6]{Demailly1992}. Its replacement of a potential near the
origin by a regularized maximum has a close analogue in Witt Nystr\"om
\cite[Proposition~3.3]{WittNystrom2018}. Those results concern projective
K\"ahler geometry; here the replacement is proved directly for the
specified form on a noncompact sublevel domain. Moser's deformation
method \cite[Section~4]{Moser1965}, with compact support verified below,
then transfers the ball to the original form. The radial coordinate
change is a special case of the K\"ahler symplectic-coordinate formula of
Loi and Zuddas \cite[proof of Theorem~1.1, equation~(23)]{LoiZuddas2008};
its pullback is also verified directly.


\subsection*{The construction}

The lower bound comes from comparing two scales. For a fixed body $K$
given by finitely many symmetric strips and a large integer $k$, a
holomorphic construction produces balls of every capacity below $\pi k$ in a suitable symplectic
domain. We realize that domain in
$\interior K\times S_k\interior K^\circ$, where the momentum scale
satisfies $S_k=(1+o(1))\pi k/4$ as $k\to\infty$.
Dividing momentum by $S_k$ and compensating with a dilation of the
source ball gives each prescribed capacity below $4$ once $k$ is
sufficiently large.
Two independent analytic ingredients supply these scales.

The first turns holomorphic vanishing into a ball
(Theorem~\ref{thm:holomorphic-ball}). Let $f=(f_1,\ldots,f_m)$ be holomorphic on an open set
$\Omega\subset\C^n$ containing $0$, have the unique common zero $0$,
and satisfy
$|f(z)|=O(|z|^k)$ there. The domain $\{|f|^2<1\}$, equipped with a symplectic form
constructed from the potential $|z|^2+|f(z)|^2$, contains balls of every
capacity strictly below $\pi k$, provided the closed sublevels of
$|f|^2$ below $1$ are compact in $\Omega$. The order of vanishing controls the
slope of a logarithmic potential near zero. Replacing its singularity
by a smooth radial potential exposes a ball, and a compactly supported
Moser deformation transports that ball to the original form.
Compact sublevels ensure the deformation is supported away from the
boundary; no regularity of the level hypersurfaces is required.

The second ingredient controls the momentum scale
(Proposition~\ref{prop:uniform-slice}). A conformal map $F$ takes the unit
disk $\D=\{w\in\C:|w|<1\}$ to a bounded domain $D$ with tips at $\pm i$
and centered horizontal slices at every height $-1<t<1$.
Write $g=F^{-1}$. The horizontal primitives of the densities
$|(g^k)'|^2$, taken from the imaginary axis, have absolute value at most
$(\pi k/4)|g|^{2k}+o(k)$, uniformly on all of $D$. Uniformity near the tips is essential:
the slice height may approach $1$ or $-1$ as $k$ increases.
Section~\ref{sec:planar} proves this estimate locally, including the
moving-endpoint bound needed at those tips.

To combine the ingredients, write
$K=\{x:|\langle b_j,x\rangle|\le1,\ 1\le j\le m\}$, where the real
vectors $b_j$ span $\R^n$. Extend these pairings complex-linearly and
set $f_j(z)=g(\langle b_j,z\rangle)^k$ on the domain where every
$\langle b_j,z\rangle$ lies in $D$. These functions have the required
unique common zero and compact sublevels. The phase-space realization
sends the imaginary part of $z$ to position and uses the horizontal
primitives for momentum. Their monotonicity makes the map globally
injective, and their uniform estimate bounds the momentum in the polar
body. Section~\ref{sec:realization} carries out this construction.
Finally, approximation in the polar body extends the lower bound to
arbitrary $K$. A supporting-cylinder argument and nonsqueezing give
the matching upper bound.

\subsection*{Organization}

Section~\ref{sec:ball} proves the holomorphic ball principle.
Section~\ref{sec:planar} constructs the planar coordinate and proves the
uniform slice estimate. Section~\ref{sec:realization} treats finite strip
bodies, and Section~\ref{sec:approximation} passes to all convex bodies,
proves the matching upper bound, and derives the volume-product and
packing consequences.
Balls and symplectic domains are open unless explicitly stated otherwise.
